% 依赖项（必须放在主文档导言区，即 \begin{document} 之前）：
% \usepackage{tabularray}
% \UseTblrLibrary{varwidth}
% \usepackage{enumitem}
\begingroup
% \nopagebreak[4]
\enlargethispage{1cm}
\footnotesize
\centering

\DefTblrTemplate{conthead-text}{default}{（续表）}
\DefTblrTemplate{contfoot-text}{default}{续下页}

\begin{longtblr}[
  caption = {北美AI智能体领域出口管制政策汇总},
  label   = {tab:north_america_ai_export_control},
]{
  width   = \linewidth,
  colspec = {
    |Q[wd=0.8cm,l,m]
    |Q[wd=1.5cm,l,m]
    |Q[wd=2.5cm,l,m]
    |Q[wd=2cm,l,m]
    |X[1.00,l,m]|
  },
  cells   = {l,m},
  colsep  = 3pt,
  row{1}  = {font=\bfseries},
  rowhead = 1,
  rowsep  = 5pt,
  measure = vbox,
  stretch = -1,
}
\hline
国家 & 主管部门 & 政策法规/管制工具 & 发布/生效日期 & 核心内容 \\
\hline
美国
% & \begin{itemize}[tabitem]
%     \item 国会
%     \item 工业和安全局(BIS)
% \end{itemize}
% & 《出口管制改革法案》ECRA+《出口管理条例》EAR~\cite{ECRA,EAR}
% & ECRA 2018；EAR持续修订
% & \begin{itemize}[tabitem]
%     \item ECRA确立美国对军民两用、敏感和新兴基础技术实施出口管制的法定基础；
%     \item EAR通过CCL/ECCN、许可要求、目的地、最终用户和最终用途规则实施具体管制；
%     \item BIS可通过实体清单、未经验证清单、军事最终用户规则和特定 FDPR 扩展执法覆盖范围。
% \end{itemize} \\
% \cline{2-5}
% & \begin{itemize}[tabitem]
%     \item 商务部BIS
% \end{itemize}
& \SetCell[r=3]{l,m} 工业和安全局(BIS)
& 2025先进AI技术全域扩散管制最终规则~\cite{AI_Diffusion_Rule}
& 2025.01.15生效（2025.05.13被废除）~\cite{Rescission_of_Biden-Era_AI_Diffusion_Rule}
& \begin{itemize}[leftmargin=*,topsep=0pt,partopsep=0pt,itemsep=1pt,parsep=0pt]
    \item 建立面向先进AI芯片和部分闭源AI模型权重的全球许可框架；
    \item 通过国家分组、受信终端用户和许可例外/授权机制区分盟友、受信主体和受限目的地；
    \item 旨在降低先进计算芯片和部分闭源模型权重被转移、规避管制或用于危害美国国家安全与外交政策利益用途的风险。
\end{itemize} \\
\cline{3-5}
% & \begin{itemize}[tabitem]
%     \item 商务部BIS
% \end{itemize}
&
& 2026对华先进芯片审查修订规则~\cite{Revises_License_Review_Policy_for_Semiconductors_Exported_to_China}
& 2026.01.15生效
& \begin{itemize}[leftmargin=*,topsep=0pt,partopsep=0pt,itemsep=1pt,parsep=0pt]
    \item 对NVIDIA H200、AMD MI325X及类似芯片的对华出口许可申请，在满足规定安全条件后实行逐案审查；
    \item 申请人须证明相关出口不会减少目前可供美国客户使用的全球半导体产能；
    \item 中国购买方须建立出口合规和客户筛查程序，相关产品还须在美国接受独立第三方性能与安全测试。
\end{itemize} \\
\cline{3-5}
% & \begin{itemize}[tabitem]
%     \item 商务部BIS
% \end{itemize}
&
& 实体清单/UVL（AI专项增补）~\cite{Entity_List}
& 2024--2026动态扩容
& \begin{itemize}[leftmargin=*,topsep=0pt,partopsep=0pt,itemsep=1pt,parsep=0pt]
    \item 实体清单对被列名主体设置额外许可要求和许可审查政策；
    \item 未经验证清单用于标记 BIS 无法核验最终用途/最终用户的交易对象，并增加出口商尽调义务；
    \item AI 半导体、模型或技术相关限制应回到具体 EAR 条款、实体清单条目和最终用途规则逐条判断。
\end{itemize} \\
\cline{2-5}
& \begin{itemize}[leftmargin=*,topsep=0pt,partopsep=0pt,itemsep=1pt,parsep=0pt]
    \item 国际贸易管理局(ITA)
    \item BIS
\end{itemize}
& 《美国人工智能出口计划》配套规则~\cite{USA_AI_Exports_Program}
& 2026.03落地
& \begin{itemize}[leftmargin=*,topsep=0pt,partopsep=0pt,itemsep=1pt,parsep=0pt]
    \item 面向海外推广由AI优化计算硬件、数据中心存储、模型、网络安全措施和行业应用组成的全栈AI技术包；
    \item 设置预设联合体和按需联合体两种模式，分别服务持续性全球部署和具体项目机会；
    \item 由商务部长会同国务院及相关部门遴选项目，获批的全栈AI技术包向可信外国买方提供。
\end{itemize} \\
\hline
加拿大
% & \begin{itemize}[tabitem]
%     \item 全球事务部 GAC（出口管制司）
% \end{itemize}
& 全球事务部(GAC)
& 《进出口许可法》EIPA+《出口管制清单》ECL（第 5506 项两用战略技术）~\cite{EIPA,ECL_Guide}
& EIPA长期生效；ECL持续年度修订
& \begin{itemize}[leftmargin=*,topsep=0pt,partopsep=0pt,itemsep=1pt,parsep=0pt]
    \item EIPA授权加拿大对特定货物、技术和目的地实施出口/进口许可管理；
    \item ECL列明受控货物与技术类别，出口商需按清单、目的地、最终用户和最终用途判断许可义务；
    \item AI相关软硬件只有在落入清单条目或其他制裁/最终用途规则时，才构成受控出口。
\end{itemize} \\
\cline{3-5}
% % & \begin{itemize}[tabitem]
% %     \item 全球事务部GAC
% % \end{itemize}
% &
% & 2024先进半导体与AI算力管制修订令SOR/2024-112 ~\cite{SOR2024-112}
% & 2024.07.20生效
% & \begin{itemize}[tabitem]
%     \item 修订加拿大Export Control List中的受控物项，以对接国际出口管制义务和战略技术变化；
%     \item 涉及半导体、先进制造或相关技术的控制应以清单条目和技术参数为准；
%     \item 许可审查聚焦物项分类、目的地、最终用户、最终用途和转出口风险，而非单独的“智能体训练”标签。
% \end{itemize} \\
% \cline{3-5}
% % & \begin{itemize}[tabitem]
% %     \item 全球事务部GAC
% % \end{itemize}
&
& 2025 战略技术单边管制修正案DORS/2025-89 ~\cite{DORS2025-89}
& 2025.04.25 生效
& \begin{itemize}[leftmargin=*,topsep=0pt,partopsep=0pt,itemsep=1pt,parsep=0pt]
    \item 对Export Control List进行修订，新增或调整特定战略货物、软件和技术控制；
    \item 出口许可义务取决于清单条目、技术参数、目的地和最终用途，而不是“AI数据集/框架”的泛化名称；
    \item 企业应进行物项分类、最终用户尽调和转出口风险审查。
\end{itemize} \\
\cline{3-5}
% & \begin{itemize}[tabitem]
%     \item 全球事务部GAC
% \end{itemize}
&
& 2025 新版《出口管制清单操作指南》~\cite{ECL_Guide_2025}
& 2025.07.01生效
& \begin{itemize}[leftmargin=*,topsep=0pt,partopsep=0pt,itemsep=1pt,parsep=0pt]
    \item 指导出口商理解Export Control List的分类结构、条目说明和技术参数；
    \item 许可判断应基于物项分类、目的地、最终用户、最终用途和适用豁免；
    \item 对AI或云访问的限制应另查具体清单条目、制裁措施或许可政策，不能由指南本身推导。
\end{itemize} \\
\cline{3-5}
% & \begin{itemize}[tabitem]
%     \item 全球事务部GAC
% \end{itemize}
&
& 《敏感技术研究领域清单》——人工智能与大数据技术~\cite{STL}
& 2024--2026持续扩容更新
& \begin{itemize}[leftmargin=*,topsep=0pt,partopsep=0pt,itemsep=1pt,parsep=0pt]
    \item 将人工智能与大数据技术列为加拿大十一类敏感技术领域之一，视为支撑多领域发展的跨领域基础技术；
    \item AI敏感技术覆盖数据科学、数字孪生、生成式AI、大语言模型、全谱系机器学习技术；
    \item 清单用于外资审查、出口管制、科研安全风险识别，最终风险判定需结合技术转移场景综合研判。
\end{itemize} \\
\hline
\end{longtblr}
\endgroup
